% GENERATED FILE. Edit canonical lessons, then run npm run generate:books.
% canonical-source-hash: fnv1a64:bdacad90c36fdc9a
% canonical-lessons: ZH-C399-yaoqiu, ZH-C399-zhiyou, ZH-W399-zhi-only, ZH-C399-zhiyao, ZH-C399-manzu, ZH-C399-chengjia

\chapter{Only If}
\label{ch:zh-only-if}
\hlchaptermodality{399}

\begin{chapteropening}
\textbf{By the end of this chapter:} \emph{I can talk about requirements, the only way something happens, what is enough, meeting a demand and settling down.}

Complete the last lesson of chapter 399: I can talk about requirements, the only way something happens, what is enough, meeting a demand and settling down.
\end{chapteropening}

\section[yāoqiú]{yāoqiú — to require; to demand (of someone)}

\label{lesson:ZH-C399-yaoqiu}



\begin{quote}
Before the new one: say the Mandarin for ready-made, then the Mandarin for to hinge on.
\end{quote}

\subsection*{You'll want to know: yāoqiú}
\emph{yāoqiú} — \textquotedblleft{}to require; to demand (of someone)\textquotedblright{}.

\begin{quote}
\emph{lǎoshī yāoqiú wǒmen měitiān shuō hànyǔ} — The teacher asks us to speak Chinese every day.
\end{quote}

Its characters come later; for now it is a word for the ear.

\begin{cousinweb}
\emph{yāoqiú} is built from \emph{yāo} 'to press' and \emph{qiú} 'to seek', both inherited from Middle Chinese, \emph{'jiew} and \emph{gjuw} (Guangyun, in Baxter's transcription), which glosses the first 'to force, in common speech' and the second 'to look for'.
\end{cousinweb}

\begin{grammarlens}[title={What you've built}]
One more word for saying what would happen.
\end{grammarlens}

\subsection*{Your turn}
\begin{itemize}
\raggedright
  \item \emph{Say it:} \emph{yāoqiú}
  \item \emph{Say it:} \emph{yāoqiú}, once more
  \item \emph{Say it:} \emph{lǎoshī yāoqiú wǒmen měitiān shuō hànyǔ} — The teacher asks us to speak Chinese every day.
  \item \emph{Recall:} say the Mandarin for ready-made, then the Mandarin for to hinge on, then say \emph{yāoqiú} again
\end{itemize}

\begin{tcolorbox}[breakable,colback=teal!4,colframe=teal!35!black,title={Before you move on}]
What does yāoqiú mean? (\textquotedblleft{}To require; to demand (of someone)\textquotedblright{}.) Say it once more.
\end{tcolorbox}
\section[\texorpdfstring{\zh{只有} (zhǐyǒu)}{只有 (zhǐyǒu)}]{\zh{只有} (zhǐyǒu) — only if (with cái after it)}

\label{lesson:ZH-C399-zhiyou}



\begin{quote}
Before the new one: say the Mandarin for to hinge on, then the Mandarin for to demand (of someone).
\end{quote}

\subsection*{You'll want to know: \zh{只有}}
\textbf{\zh{只有}} — \emph{zhǐyǒu} — \textquotedblleft{}only if (with cái after it)\textquotedblright{}.

\begin{quote}
\textbf{\zh{只有你去他才会去}} \emph{zhǐyǒu nǐ qù tā cái huì qù} — He'll only go if you go.
\end{quote}

The new character, in pieces: \zh{只} — five strokes: \zh{口} on top, and \zh{八} under it. The next lesson puts it in your hand.

Two third tones meet in this word: zhǐ is said with a rising tone before yǒu, and the spelling keeps both marks.

\begin{cousinweb}
\textbf{\zh{只有}} (\emph{zhǐyǒu}) is built from \emph{zhǐ} 'only' and \zh{有} 'to have'. The second is inherited from Middle Chinese \emph{hjuwX} (Guangyun, in Baxter's transcription); the dictionary lists the first, \emph{tsyeX} there, as a particle.
\end{cousinweb}

\begin{grammarlens}[title={What you've built}]
One more word for saying what would happen.
\end{grammarlens}

\subsection*{Your turn}
\begin{itemize}
\raggedright
  \item \emph{Say it:} \emph{zhǐyǒu}
  \item \emph{Say it:} \emph{zhǐyǒu}, once more
  \item \emph{Say it:} \emph{zhǐyǒu nǐ qù tā cái huì qù} — He'll only go if you go.
  \item \emph{Recall:} say the Mandarin for to hinge on, then the Mandarin for to demand (of someone), then say \emph{zhǐyǒu} again
\end{itemize}

\begin{tcolorbox}[breakable,colback=teal!4,colframe=teal!35!black,title={Before you move on}]
What does \zh{只有} mean? (\textquotedblleft{}Only if (with cái after it)\textquotedblright{}.) Say it once more.
\end{tcolorbox}
\section[\texorpdfstring{\zh{只} (zhǐ)}{只 (zhǐ)}]{\zh{只} — \zh{口} over \zh{八}}

\label{lesson:ZH-W399-zhi-only}



\begin{quote}
The word first, then the character inside it.

\begin{itemize}
\raggedright
  \item \emph{Recall:} say \emph{zhǐyǒu} — \textbf{R1}, one lesson back, before it decays
\end{itemize}

\begin{quote}
\zh{口} \zh{八}
\end{quote}

\emph{Read it:} each one aloud, with nothing said first — \emph{kǒu}, \emph{bā}
\end{quote}

\subsection*{Script — \zh{口} over \zh{八}}
\hlblockfigure{\detokenize{figures/ZH-W399-zhi-only-filmstrip.pdf}}{How \zh{只} is written, stroke by stroke}

\begin{quote}
\zh{只}
\end{quote}

\textbf{Two pieces, top before bottom.}

\textbf{\zh{口} on top, three strokes,} as you wrote it: a short vertical, a stroke that runs right and turns down, then the bar that closes the foot.

\textbf{\zh{八} under it, two strokes:} a stroke falling to the left, then a dot falling to the right.

3 + 2 = five.

\subsection*{Your turn}
\begin{itemize}
\raggedright
  \item \emph{Write it:} \zh{只}, all five, with the model in front of you
  \item \emph{Cover:} write it again from memory, and count the pieces, not the strokes
\end{itemize}

\begin{tcolorbox}[breakable,colback=teal!4,colframe=teal!35!black,title={Before you move on}]
How many strokes in \zh{只}? (\textbf{Five}.) What is on top? (\textbf{\zh{口}}.) What is under it? (\textbf{\zh{八}}.)

Source: \href{https://raw.githubusercontent.com/chanind/hanzi-writer-data/68d10a4b21150cae5e1ebbd223eed289cf32d90c/data/\%E5\%8F\%AA.json}{Hanzi Writer Data \zh{只}, pinned stroke paths and medians}.
\end{tcolorbox}
\section[\texorpdfstring{\zh{只要} (zhǐyào)}{只要 (zhǐyào)}]{\zh{只要} (zhǐyào) — as long as; provided (that)}

\label{lesson:ZH-C399-zhiyao}



\begin{quote}
Before the new one: say the Mandarin for to demand (of someone), then the Mandarin for only if.
\end{quote}

\subsection*{You'll want to know: \zh{只要}}
\textbf{\zh{只要}} — \emph{zhǐyào} — \textquotedblleft{}as long as; provided (that)\textquotedblright{}.

\begin{quote}
\textbf{\zh{只要不下雨我们就去}} \emph{zhǐyào bú xiàyǔ wǒmen jiù qù} — As long as it doesn't rain, we'll go.
\end{quote}

\begin{cousinweb}
\textbf{\zh{只要}} (\emph{zhǐyào}) is built from \zh{只} 'only' and \zh{要} 'to need'. The second is inherited from Middle Chinese \emph{'jiewH} (Guangyun, in Baxter's transcription), glossed 'to bind by agreement'; the dictionary lists the first, \emph{tsyeX} there, as a particle.
\end{cousinweb}

\begin{grammarlens}[title={What you've built}]
One more word for saying what would happen.
\end{grammarlens}

\subsection*{Your turn}
\begin{itemize}
\raggedright
  \item \emph{Say it:} \emph{zhǐyào}
  \item \emph{Say it:} \emph{zhǐyào}, once more
  \item \emph{Say it:} \emph{zhǐyào bú xiàyǔ wǒmen jiù qù} — As long as it doesn't rain, we'll go.
  \item \emph{Recall:} say the Mandarin for to demand (of someone), then the Mandarin for only if, then say \emph{zhǐyào} again
\end{itemize}

\begin{tcolorbox}[breakable,colback=teal!4,colframe=teal!35!black,title={Before you move on}]
What does \zh{只要} mean? (\textquotedblleft{}As long as; provided (that)\textquotedblright{}.) Say it once more.
\end{tcolorbox}
\section[mǎnzú]{mǎnzú — to meet (a requirement); to satisfy}

\label{lesson:ZH-C399-manzu}



\begin{quote}
Before the new one: say the Mandarin for only if, then the Mandarin for as long as.
\end{quote}

\subsection*{You'll want to know: mǎnzú}
\emph{mǎnzú} — \textquotedblleft{}to meet (a requirement); to satisfy\textquotedblright{}.

\begin{quote}
\emph{zhǐyào mǎnzú zhèxiē yāoqiú nǐ jiù néng qù} — As long as you meet these requirements, you can go.
\end{quote}

Its characters come later; for now it is a word for the ear.

\begin{cousinweb}
\emph{mǎnzú} is built from \emph{mǎn} 'full' and \emph{zú} 'enough', both inherited from Middle Chinese, \emph{manX} and \emph{tsjowk} (Guangyun, in Baxter's transcription), which glosses the first 'brimming' and the second 'the foot' and 'full'.
\end{cousinweb}

\begin{grammarlens}[title={What you've built}]
One more word for saying what would happen.
\end{grammarlens}

\subsection*{Your turn}
\begin{itemize}
\raggedright
  \item \emph{Say it:} \emph{mǎnzú}
  \item \emph{Say it:} \emph{mǎnzú}, once more
  \item \emph{Say it:} \emph{zhǐyào mǎnzú zhèxiē yāoqiú nǐ jiù néng qù} — As long as you meet these requirements, you can go.
  \item \emph{Recall:} say the Mandarin for only if, then the Mandarin for as long as, then say \emph{mǎnzú} again
\end{itemize}

\begin{tcolorbox}[breakable,colback=teal!4,colframe=teal!35!black,title={Before you move on}]
What does mǎnzú mean? (\textquotedblleft{}To meet (a requirement); to satisfy\textquotedblright{}.) Say it once more.
\end{tcolorbox}
\section[\texorpdfstring{\zh{成家} (chéngjiā)}{成家 (chéngjiā)}]{\zh{成家} (chéngjiā) — to marry and settle down}

\label{lesson:ZH-C399-chengjia}



\begin{quote}
Before the new one: say the Mandarin for as long as, then the Mandarin for to meet (a requirement).
\end{quote}

\subsection*{You'll want to know: \zh{成家}}
\textbf{\zh{成家}} — \emph{chéngjiā} — \textquotedblleft{}to marry and settle down\textquotedblright{}.

\begin{quote}
\textbf{\zh{只要找到好工作他就想成家}} \emph{zhǐyào zhǎodào hǎo gōngzuò tā jiù xiǎng chéngjiā} — As long as he finds a good job, he wants to marry and settle down.
\end{quote}

\begin{cousinweb}
\textbf{\zh{成家}} (\emph{chéngjiā}) is built from \zh{成} 'to complete' and \zh{家} 'a home', both inherited from Middle Chinese, \emph{dzyeng} and \emph{kae} (Guangyun, in Baxter's transcription), which glosses the second 'a dwelling'.
\end{cousinweb}

\begin{grammarlens}[title={What you've built}]
One more word for saying what would happen.
\end{grammarlens}

\subsection*{Your turn}
\begin{itemize}
\raggedright
  \item \emph{Say it:} \emph{chéngjiā}
  \item \emph{Say it:} \emph{chéngjiā}, once more
  \item \emph{Say it:} \emph{zhǐyào zhǎodào hǎo gōngzuò tā jiù xiǎng chéngjiā} — As long as he finds a good job, he wants to marry and settle down.
  \item \emph{Recall:} say the Mandarin for as long as, then the Mandarin for to meet (a requirement), then say \emph{chéngjiā} again
\end{itemize}

\begin{tcolorbox}[breakable,colback=teal!4,colframe=teal!35!black,title={Before you move on}]
What does \zh{成家} mean? (\textquotedblleft{}To marry and settle down\textquotedblright{}.) Say it once more.
\end{tcolorbox}
